% AUTO-GENERATED by tools/usc_extract.py — do not edit by hand
\begingroup\small
\begin{historicalsource}[{\textbf{Primary text} --- Title~18, \S~1589 (Forced labor) (OLRC via \texttt{nickvido/us-code}, tag \texttt{annual/2025})}]
\textit{Source:} \url{https://uscode.house.gov/view.xhtml?req=granuleid:USC-prelim-title18&num=0&edition=prelim}\par\medskip\textbf{§ 1589. Forced labor}\\
\textbf{(a)} Whoever knowingly provides or obtains the labor or services of a person by any one of, or by any combination of, the following means—

(1) by means of force, threats of force, physical restraint, or threats of physical restraint to that person or another person;

(2) by means of serious harm or threats of serious harm to that person or another person;

(3) by means of the abuse or threatened abuse of law or legal process; or

(4) by means of any scheme, plan, or pattern intended to cause the person to believe that, if that person did not perform such labor or services, that person or another person would suffer serious harm or physical restraint,
shall be punished as provided under subsection (d).

\textbf{(b)} Whoever knowingly benefits, financially or by receiving anything of value, from participation in a venture which has engaged in the providing or obtaining of labor or services by any of the means described in subsection (a), knowing or in reckless disregard of the fact that the venture has engaged in the providing or obtaining of labor or services by any of such means, shall be punished as provided in subsection (d).

\textbf{(c)} In this section:

(1) The term “abuse or threatened abuse of law or legal process” means the use or threatened use of a law or legal process, whether administrative, civil, or criminal, in any manner or for any purpose for which the law was not designed, in order to exert pressure on another person to cause that person to take some action or refrain from taking some action.

(2) The term “serious harm” means any harm, whether physical or nonphysical, including psychological, financial, or reputational harm, that is sufficiently serious, under all the surrounding circumstances, to compel a reasonable person of the same background and in the same circumstances to perform or to continue performing labor or services in order to avoid incurring that harm.

\textbf{(d)} Whoever violates this section shall be fined under this title, imprisoned not more than 20 years, or both. If death results from a violation of this section, or if the violation includes kidnaping, an attempt to kidnap, aggravated sexual abuse, or an attempt to kill, the defendant shall be fined under this title, imprisoned for any term of years or life, or both.
\end{historicalsource}
\endgroup
